\section{Design}
\label{sec:design}

\sysname's design centres on three primitives: the \captoken, the
\texttt{CapabilityManager} that mints and revokes them, and the
\capgate that fronts every IPC operation.

\subsection{Token format}
\label{sec:design:token}

A \captoken is a 192-bit value type (\Cref{fig:token}). Tokens are
passed by value, copied freely, and contain no pointers. The
authoritative state lives in the \texttt{CapabilityManager}'s pool;
the token is the keying material that names a slot.

\begin{figure}[t]
\centering\small
\begin{tabular}{|c|c|}
\hline
\textbf{Field} & \textbf{Bits} \\
\hline
\texttt{m\_arrUId[0]} : slot index   & 32 \\
\texttt{m\_arrUId[0]} : random       & 32 \\
\texttt{m\_arrUId[1]} : random       & 64 \\
\texttt{m\_ePermissions}             & 64 \\
\texttt{m\_uEpoch}                   & 64 \\
\texttt{m\_uOwnerId}                 & 64 \\
\texttt{m\_arrUParentId}             & 128 \\
\hline
\end{tabular}
\caption{Capability-token wire layout. The \texttt{m\_arrUId[0]}
high half is the slot-index hint that powers the $O(1)$ validate
path; the remaining 96 bits are RDRAND randomness.}
\label{fig:token}
\end{figure}

The 128-bit identifier (\texttt{m\_arrUId}) splits as 32 bits of slot
index plus 96 bits of RDRAND randomness. The split is the central
design choice and is justified in \Cref{sec:design:fastpath}.

\subsection{Permission monotonicity}
\label{sec:design:mono}

Children derived from a parent token may carry a permission set that
is a strict subset of the parent's. \texttt{CapabilityManager::derive}
enforces this invariant; an attempt to derive with a non-subset
permission mask returns \texttt{Result::error}. The CBMC bounded-model
proof (\Cref{sec:impl:cbmc}) checks this property over all possible
parent / child permission pairs in the supported bitmask.

The motivation is the standard capability discipline: once a
permission has been removed in the derivation chain, it cannot be
recovered without going back to a higher-permission ancestor.
\sysname does not currently support \textit{regrant} (re-issuing a
permission that has been revoked); this is by design.

\subsection{The slot-indexed validate fast path}
\label{sec:design:fastpath}

The most consequential design decision is the layout of the
\texttt{validate} path.

\para{Why the natural design is too slow}
A typical capability validate path looks roughly like:
\begin{enumerate}
\itemsep0pt
\item Hash the 128-bit token identifier to an index in a side table.
\item Load the bucket; walk a chain.
\item On match, check epoch, check permission, check
\texttt{m\_uRevokedAtEpoch}.
\end{enumerate}
On a 4096-token pool, even with a perfect hash this costs roughly one
load (the bucket head) plus one constant-time compare; on a large
pool with collision chains, it can cost two or more cache misses
before the compare even fires. On the IPC fastpath, where the
producer and consumer each call \texttt{validate} once per message,
this dwarfs the rest of the cost.

\para{Our approach}
We embed the slot index inside the token's identifier itself:

\begin{equation*}
\texttt{m\_arrUId[0]} \;=\; (\text{slot}_{32} \ll 32) \;|\;
                            \text{rand}_{32}
\end{equation*}

\noindent
\texttt{validate} then runs:

\begin{verbatim}
if (!token.isValid())    return false;
// epoch sanity: reject tokens claiming
// an epoch that hasn't happened yet
if (token.epoch > current_epoch)
    return false;
if ((token.perms & required) != required)
    return false;
slot = token.uid[0] >> 32;
if (slot >= max_tokens)  return false;
TokenSlot& s = pool[slot];  // 1 load
if (!s.active.load(acquire))    return false;
if (s.revoked_at.load(acquire)) return false;
return asm_ct_compare(s.token.uid,
                      token.uid, 16) == 0;
\end{verbatim}

\noindent
The register-only checks (epoch sanity, permission subset, slot
bounds) run before the single pool-slot load; the constant-time
compare on the full 128-bit identifier is the final arbiter. This
listing mirrors the shipped \texttt{CapabilityManager::validate}
check-for-check, in order.

\para{Why this is safe}
The slot index is treated as a \textit{hint}, not a credential. An
attacker who forges a slot index gets a deterministic indirection
into the pool; the constant-time compare on the full 128 bits then
has to match. The forgery space remains $2^{96}$. The bounds check
on the slot makes out-of-range indices reject without an
out-of-bounds load (\Cref{sec:eval:correctness}, T-OOB).

\para{Why this is fast}
The validate path's data-dependent work is a single shift, a single
load (\texttt{pool[slot]}), and a single \texttt{ct\_compare} on
16~bytes. The cache load is the only memory access; the rest is in
registers. \Cref{fig:figure-2} confirms the cost is flat across pool
sizes from 16 to 4090 active tokens.

\para{Constant-time compare}
The 128-bit identifier compare uses the M-21
\texttt{asm\_ct\_compare} primitive, which fans the bytes through a
running OR and returns a single-bit verdict. This is necessary
because a non-constant-time compare would leak which byte differs
first, allowing a chosen-token attack to recover slots that have
already been allocated.

\subsection{Cascading revocation}
\label{sec:design:revoke}

\texttt{revoke(t)} locates $t$'s slot in $O(1)$ via the same
slot-index trick as validate, stamps the slot's
\texttt{m\_uRevokedAtEpoch}, publishes deactivation with a release
store, and securely wipes the token bytes
(\texttt{asm\_secure\_wipe}). Any \texttt{validate()} that runs after
the release store rejects on the slot flags; a racing
\texttt{validate()} that read \texttt{active==true} just before
deactivation still rejects, because the constant-time compare runs
against wiped bytes. Each revoke transaction also increments a
global epoch counter, which serves two roles: a sanity check
(validate rejects tokens claiming a future epoch) and a shared
timestamp for audit coherence---every token revoked in one cascade
carries the same epoch stamp.

The cascade then revokes all live descendants of $t$. Tokens carry a
parent \emph{backpointer} (\texttt{m\_arrUParentId}) but no forward
child list---a deliberate choice that keeps tokens fixed-size value
types with no heap state---so the cascade discovers children by
sweeping the bounded pool once per BFS layer: $O(N)$ per visited
node with $N \le 4096$. We report this honestly rather than claiming
$O(d)$: the cascade is bookkeeping that runs under the manager's
spinlock, \emph{off} the message fastpath. The property that matters
for the security argument---and the one we measure in
\Cref{sec:eval:revoke}---is \emph{time-to-effect}: the interval from
the \texttt{revoke()} call to the first rejected message, which is
independent of pool occupancy because rejection rides on the
$O(1)$ slot-flag check inside validate.

A revoked slot is freed for reuse on the next \texttt{create()};
because the wipe destroys the 96 random bits, a stale token can
never match a recycled slot.

\subsection{The capability gate on the IPC fastpath}
\label{sec:design:gate}

The \capgate binds a \texttt{CapabilityManager*} and a
\texttt{CapabilityToken} to one endpoint of an IPC channel
(\Cref{fig:gate}). A typical \sysname channel has two gates: one
stamped \IPCSEND held by the producer endpoint, one stamped \IPCRECV
held by the consumer endpoint. Each \texttt{writeGated()} call
invokes \texttt{validate()} \textit{before} attempting to claim ring
space; each \texttt{readGated()} invokes \texttt{validate()} before
attempting to consume.

\begin{figure}[t]
\centering
\includegraphics[width=0.95\columnwidth]{figures/gate_path}
\caption{The capability gate on the IPC fastpath. The
\texttt{CapabilityGate} sits in front of the existing
\texttt{RingBuffer::write/read} machinery; on validate failure no
ring state is mutated.}
\label{fig:gate}
\end{figure}

\para{The denied-request invariant}
A denied request must leave the ring's
\texttt{m\_uWriteClaimIndex} and \texttt{m\_uWriteIndex} unchanged.
Without this property, a revoked principal could DOS a channel by
burning claim space without ever committing. \Cref{sec:eval:correctness}
T-RingFP exercises this invariant directly; the test asserts that
neither index advances after any of the three deny modes (unbound
gate, missing permission, revoked token).

\para{The blocking-read invariant}
\texttt{readWaitGated()} re-validates after every futex wake-up. A
token can be revoked while a reader is parked on the futex; if we
delivered a message past the revoke epoch on a stale wake, the gate
would have a window. Re-validation closes the window at one extra
\texttt{validate()} call per wake (sub-100\,ns).
